\documentclass[11pt]{article}
\usepackage{amsmath,amssymb,amsthm,mathtools}
\usepackage{enumitem}
\usepackage[margin=1in]{geometry}
\usepackage{hyperref}

\title{Optimal Constant Step Size for Population GD in Gaussian-Sketched Random Features\\
via Weighted Spectral Regularity on Relevant Scales\\
\small (with curvature at the optimum)}
\author{}
\date{}

\newtheorem{theorem}{Theorem}
\newtheorem{lemma}{Lemma}
\newtheorem{assumption}{Assumption}
\newtheorem{proposition}{Proposition}
\newtheorem{remark}{Remark}

\begin{document}
\maketitle

\paragraph{Scope (what is proved vs.\ what is imported).}
The deterministic core of this note is fully self-contained at the level of
\emph{finite-dimensional} quadratic objectives: for any SPD $H\in\mathbb R^{d\times d}$,
any $b\in\mathbb R^d$, and any horizon $T\in\mathbb N$, we have an exact spectral
representation of $L_T(\eta;H,b)$ and its derivatives.

The asymptotic expansions in Section~\ref{sec:det} are \emph{conditional} on a
$b$-weighted spectral regularity assumption on the window $\mu\sim 1/T$
(Assumption~\ref{ass:wRV}), and should be interpreted \emph{uniformly along a family of
finite-dimensional instances} $(H_N,b_N)$ with horizons $T=T(N)\to\infty$ (e.g.\ the RF
setting with $d=N$). In particular, for a single fixed matrix $H$ in fixed dimension,
Assumption~\ref{ass:wRV} can only hold up to horizons $T\lesssim 1/\mu_d(H)$
(see Remark~\ref{rem:finite-d}).

A second (and independent) uniformity issue concerns the contribution of the
``head'' modes $\{\mu_j>\mu_{\mathrm{cut}}\}$: for a fixed $H$ this head is exponentially
small in $T$, but to obtain \emph{uniform} asymptotics along $(H_N,b_N)$ one needs a
uniform control on the top edge of the spectrum. In the Gaussian RF model this is
provided by a population eigengap $\lambda_1>\lambda_2$ and operator-norm concentration
(Section~\ref{sec:rf}).

\begin{itemize}[leftmargin=*]
\item[(D1)] Under the $b$-weighted spectral regularity assumption on scales $\mu\sim 1/T$
and a mild uniform top-edge separation condition (automatic in the RF model),
we prove a deterministic step-size minimizer expansion for constant-step population GD
on a quadratic objective (Theorem~\ref{thm:det}).
\item[(D2)] Under the same assumptions, we prove a deterministic \emph{curvature-at-the-optimum}
result (Theorem~\ref{thm:det-curv}), showing $L_T''(\eta_T^\star)>0$ and giving sharp
two-sided scaling for $L_T''(\eta_T^\star)$ in terms of the weighted spectral exponent.
\end{itemize}

For the Gaussian random-features (RF) model, we assume standard \emph{capacity} and
\emph{source} conditions on the population covariance operator $\Sigma$ and target
$\bar w$, and we appeal to deterministic-equivalent / local-law results to justify the
RF weighted spectral regularity assumption on the scale $\mu\sim 1/T$.

\begin{remark}[Finite dimension vs.\ asymptotics in the deterministic core]
Throughout Section~\ref{sec:det}, all identities and derivative manipulations are
carried out for finite-dimensional SPD matrices $H\in\mathbb R^{d\times d}$. However,
the asymptotic statements in Theorems~\ref{thm:det} and~\ref{thm:det-curv} are intended
to be applied along sequences of finite-dimensional instances $(H_N,b_N)$ with horizons
$T=T(N)\to\infty$ for which the weighted spectral regularity window
$\mu\in[1/T,\mu_{\mathrm{cut}}]$ remains above the spectral floor.
This is exactly the RF regime where $d=N\to\infty$ and $T\asymp N^\alpha$.
\end{remark}

% ------------------------------------------------------------------
\section{Deterministic population GD and a weighted spectral viewpoint}
\label{sec:det}

\subsection{Quadratic setup and exact loss representation}

Let $H\in\mathbb R^{d\times d}$ be symmetric positive definite (SPD). Assume $H$ has
orthonormal eigenvectors $\{v_j\}_{j=1}^d$ with eigenvalues
\[
\mu_1>\mu_2\ge \mu_3\ge \cdots \ge \mu_d>0.
\]
Fix $b\in\mathbb R^d$ and consider the quadratic objective
\[
\mathcal L(v):=\tfrac12\langle v,Hv\rangle-\langle b,v\rangle+\mathrm{const}.
\]
The unique minimizer is $v^\star=H^{-1}b$. Run population gradient descent with
constant step size $\eta>0$ and initialization $v^{(0)}=0$:
\[
v^{(t+1)}=v^{(t)}-\eta(Hv^{(t)}-b),\qquad t=0,1,2,\dots
\]
Let $e^{(t)}:=v^{(t)}-v^\star$. Then
\[
e^{(t)}=(I-\eta H)^t e^{(0)},\qquad e^{(0)}=-v^\star.
\]
Write $b_j:=\langle b,v_j\rangle$. Since $Hv^\star=b$, we have
$\mu_j\langle v^\star,v_j\rangle=b_j$, hence $\langle v^\star,v_j\rangle=b_j/\mu_j$.
Therefore the excess loss after $T$ steps is
\begin{equation}\label{eq:LT}
L_T(\eta;H,b):=\mathcal L(v^{(T)})-\mathcal L(v^\star)
=\tfrac12\sum_{j=1}^d\frac{b_j^2}{\mu_j}\,(1-\eta\mu_j)^{2T}.
\end{equation}
Since $d<\infty$, $L_T$ is a polynomial in $\eta$ and thus $C^\infty$. Differentiating
term-by-term yields
\begin{align}
L_T'(\eta;H,b)
&=-T\sum_{j=1}^d b_j^2(1-\eta\mu_j)^{2T-1},
\label{eq:LTprime}\\
L_T''(\eta;H,b)
&=
T(2T-1)\sum_{j=1}^d b_j^2\,\mu_j\,(1-\eta\mu_j)^{2T-2}.
\label{eq:LTsecond}
\end{align}

\begin{lemma}[Convexity in $\eta$]\label{lem:convex}
For every $T\ge 1$, $L_T''(\eta)\ge 0$ for all $\eta\in\mathbb R$. Hence $L_T$ is convex
on $\mathbb R$, and any interior stationary point is a global minimizer.
\end{lemma}

\begin{proof}
In \eqref{eq:LTsecond}, the exponent $2T-2$ is even and $\mu_j>0$, so each summand is
nonnegative.
\end{proof}

\begin{assumption}[Top-mode nondegeneracy]\label{ass:top}
$b_1=\langle b,v_1\rangle\neq 0$.
\end{assumption}

\begin{lemma}[Strict convexity on $(1/\mu_1,2/\mu_1)$]\label{lem:strict-convex}
Assume Assumption~\ref{ass:top} and $T\ge 2$. Then $L_T''(\eta)>0$ for every
$\eta\in(1/\mu_1,2/\mu_1)$, and in particular the minimizer $\eta_T^\star$ defined below
is unique.
\end{lemma}

\begin{proof}
For $\eta\neq 1/\mu_1$, the $j=1$ summand in \eqref{eq:LTsecond} equals
$T(2T-1)b_1^2\mu_1(1-\eta\mu_1)^{2T-2}>0$ since $b_1\neq 0$ and the exponent is even.
\end{proof}

% ------------------------------------------------------------------
\subsection{Weighted spectral CDF and the relevant scale $\mu\sim 1/T$}

Fix a cutoff $\mu_{\mathrm{cut}}\in(0,\mu_1/4]$ and define the $b$-weighted spectral
measure away from the top mode on $(0,\mu_{\mathrm{cut}}]$:
\[
\nu_b:=\sum_{\substack{2\le j\le d:\\ \mu_j\le \mu_{\mathrm{cut}}}} b_j^2\,\delta_{\mu_j},
\qquad
F_b(x):=\nu_b([0,x])=\sum_{\substack{2\le j\le d:\\ \mu_j\le \min\{x,\mu_{\mathrm{cut}}\}}} b_j^2.
\]
We assume a power law \emph{only for this weighted CDF} on the relevant window
$[1/T,\mu_{\mathrm{cut}}]$.

\begin{assumption}[Weighted regular variation on relevant scales]\label{ass:wRV}
Fix $\theta>0$ and $\mu_{\mathrm{cut}}\in(0,\mu_1/4]$. There exist constants
$c_F,C_F>0$ and $T_0\ge 1$ such that for all $T\ge T_0$,
\begin{equation}\label{eq:wRV}
c_F x^{\theta+1}\le F_b(x)\le C_F x^{\theta+1}
\qquad \forall x\in[1/T,\mu_{\mathrm{cut}}].
\end{equation}
\end{assumption}

\begin{remark}[Finite dimension and the lower end of the window]\label{rem:finite-d}
In fixed finite dimension, $F_b(x)=0$ for $x<\mu_d$ and \eqref{eq:wRV} cannot hold down to
$x\downarrow 0$. The intended regime is that $T$ does not exceed a constant multiple of
$1/\mu_d$ so that $1/T\gtrsim \mu_d$ and the window $[1/T,\mu_{\mathrm{cut}}]$ lies above the
spectral floor. This matches RF regimes with $\mu_{\min}(H_N)\asymp N^{-\alpha}$ and
$T\lesssim N^\alpha$.
\end{remark}

\begin{assumption}[Uniform top-edge separation (for uniform-in-family bounds)]\label{ass:gap}
There exists $\varepsilon_{\mathrm{gap}}\in(0,1)$ such that
\[
\mu_2 \le (1-\varepsilon_{\mathrm{gap}})\,\mu_1.
\]
\end{assumption}

% ------------------------------------------------------------------
\subsection{From $S_T(\eta)$ to a Laplace transform}

Split \eqref{eq:LTprime} into: top mode ($j=1$), tail ($\mu_j\le \mu_{\mathrm{cut}}$),
and head ($\mu_j>\mu_{\mathrm{cut}}$, excluding $j=1$):
\[
L_T'(\eta)=
\underbrace{-T b_1^2(1-\eta\mu_1)^{2T-1}}_{\text{top}}
\;-\;T S_T(\eta)\;-\;T H_T(\eta),
\]
where
\[
S_T(\eta):=\sum_{\substack{2\le j\le d:\\ \mu_j\le\mu_{\mathrm{cut}}}}
b_j^2(1-\eta\mu_j)^{2T-1},
\qquad
H_T(\eta):=\sum_{\substack{2\le j\le d:\\ \mu_j>\mu_{\mathrm{cut}}}}
b_j^2(1-\eta\mu_j)^{2T-1}.
\]

Define the Laplace transform of $F_b$:
\[
\mathcal L_{F_b}(s):=\int_{[0,\mu_{\mathrm{cut}}]}e^{-s\mu}\,dF_b(\mu)
=\sum_{\substack{2\le j\le d:\\ \mu_j\le\mu_{\mathrm{cut}}}} b_j^2 e^{-s\mu_j}.
\]

\begin{lemma}[Elementary exponential comparison]\label{lem:exp-comp}
For $x\in[0,1/2]$,
\[
e^{-x-x^2}\le 1-x\le e^{-x}.
\]
Consequently, if $\eta\mu\le 1/2$ then for all integers $T\ge 1$,
\[
e^{-4\eta T\mu}\ \le\ (1-\eta\mu)^{2T-1}\ \le\ e^{-\eta T\mu}.
\]
Moreover, if $\eta\mu\le 1/2$ and $T\ge 2$, then
\[
e^{-4\eta T\mu}\ \le\ (1-\eta\mu)^{2T-2}\ \le\ e^{-\eta T\mu}.
\]
\end{lemma}

\begin{proof}
The bounds $1-x\le e^{-x}$ and $\log(1-x)\ge -x-x^2$ for $x\in[0,1/2]$ are standard.
For the $(2T-1)$ upper bound: $(1-x)^{2T-1}\le e^{-(2T-1)x}\le e^{-Tx}$ since $2T-1\ge T$.
For the $(2T-2)$ upper bound: $(1-x)^{2T-2}\le e^{-(2T-2)x}\le e^{-Tx}$ since $2T-2\ge T$ for $T\ge 2$.

For the lower bound, using $\log(1-x)\ge -x-x^2$ and $x^2\le x$ on $[0,1/2]$,
\[
(2T-1)\log(1-x)\ge -(2T-1)(x+x^2)\ge -2T(x+x^2)\ge -4Tx,
\]
and similarly with $2T-2$ in place of $2T-1$. Exponentiate and set $x=\eta\mu$.
\end{proof}

Let
\[
\eta_0:=\frac{2}{\mu_1}-\frac{1}{4\mu_1}=\frac{7}{4\mu_1},
\qquad
I:=[\eta_0,2/\mu_1].
\]
Then for any $\eta\in I$ and any $\mu\le\mu_{\mathrm{cut}}\le\mu_1/4$,
we have $\eta\mu\le (2/\mu_1)(\mu_1/4)=1/2$, so Lemma~\ref{lem:exp-comp} applies.

\begin{lemma}[Laplace bounds from weighted CDF on relevant scales]\label{lem:laplace}
Assume Assumption~\ref{ass:wRV}. Then there exist constants $c_L,C_L>0$ and $T_1\ge T_0$
(depending only on $\mu_1,\mu_{\mathrm{cut}},\theta,c_F,C_F$) such that for all
$T\ge T_1$ and all $s\in[c_1 T,c_2 T]$ with
\[
c_1:=\eta_0,\qquad c_2:=\frac{8}{\mu_1},
\]
we have
\[
c_L s^{-(\theta+1)}\le \mathcal L_{F_b}(s)\le C_L s^{-(\theta+1)}.
\]
\end{lemma}

\begin{proof}
Split
\[
\mathcal L_{F_b}(s)
=\int_{[0,1/T]}e^{-s\mu}\,dF_b(\mu)+\int_{(1/T,\mu_{\mathrm{cut}}]}e^{-s\mu}\,dF_b(\mu).
\]

\emph{Upper bound.}
For the first term, $e^{-s\mu}\le 1$ and $F_b$ is nondecreasing, so
\[
\int_{[0,1/T]}e^{-s\mu}\,dF_b(\mu)\le F_b(1/T)\le C_F (1/T)^{\theta+1}.
\]

For the second term, use Stieltjes integration by parts with $f(\mu)=e^{-s\mu}$:
\[
\int_{(1/T,\mu_{\mathrm{cut}}]} e^{-s\mu}\,dF_b(\mu)
=
e^{-s\mu_{\mathrm{cut}}}F_b(\mu_{\mathrm{cut}})
-e^{-s/T}F_b(1/T)
+s\int_{1/T}^{\mu_{\mathrm{cut}}} e^{-s\mu}F_b(\mu)\,d\mu.
\]
Dropping the negative boundary term $-e^{-s/T}F_b(1/T)$ and using
$F_b(\mu)\le C_F\mu^{\theta+1}$ on $[1/T,\mu_{\mathrm{cut}}]$ yields
\[
\int_{(1/T,\mu_{\mathrm{cut}}]} e^{-s\mu}\,dF_b(\mu)
\le
e^{-s\mu_{\mathrm{cut}}}F_b(\mu_{\mathrm{cut}})
+s C_F\int_{1/T}^{\mu_{\mathrm{cut}}} e^{-s\mu}\mu^{\theta+1}\,d\mu.
\]
With $t=s\mu$, the integral term is
\[
sC_F\int_{1/T}^{\mu_{\mathrm{cut}}} e^{-s\mu}\mu^{\theta+1}\,d\mu
=
C_F s^{-(\theta+1)}\int_{s/T}^{s\mu_{\mathrm{cut}}} e^{-t}t^{\theta+1}\,dt
\le C'_F s^{-(\theta+1)}.
\]
For the boundary term, since $s\ge c_1T$ on $[c_1T,c_2T]$,
\[
e^{-s\mu_{\mathrm{cut}}}F_b(\mu_{\mathrm{cut}})
\le
F_b(\mu_{\mathrm{cut}})\,e^{-c_1\mu_{\mathrm{cut}}T}
=O(e^{-cT}),
\]
which is eventually dominated by any polynomial rate in $T$ and in particular by
$s^{-(\theta+1)}$ on $s\asymp T$.
Together with $F_b(1/T)\le C_FT^{-(\theta+1)}\asymp s^{-(\theta+1)}$, this yields the upper bound.

\medskip
\emph{Lower bound.}
For $s\in[c_1T,c_2T]$ and $\mu\in[0,1/T]$, we have
$s\mu\le c_2T\cdot \frac1T = 8/\mu_1$, hence $e^{-s\mu}\ge e^{-8/\mu_1}$. Therefore
\[
\mathcal L_{F_b}(s)\ge \int_{[0,1/T]} e^{-s\mu}\,dF_b(\mu)
\ge e^{-8/\mu_1}\big(F_b(1/T)-F_b(0)\big).
\]
Since $H$ is SPD, there are no eigenvalues at $0$, so $F_b(0)=0$. Assumption~\ref{ass:wRV}
gives $F_b(1/T)\ge c_F T^{-(\theta+1)}$, and $T^{-(\theta+1)}\asymp s^{-(\theta+1)}$
on $[c_1T,c_2T]$.
\end{proof}

\begin{lemma}[Laplace bounds with one power of $\mu$]\label{lem:laplace-mu}
Assume Assumption~\ref{ass:wRV}. Then there exist constants $c_{L,1},C_{L,1}>0$
and $T_{1,1}\ge T_0$ (depending only on $\mu_1,\mu_{\mathrm{cut}},\theta,c_F,C_F$)
such that for all $T\ge T_{1,1}$ and all $s\in[c_1T,c_2T]$ (with $c_1,c_2$ as in
Lemma~\ref{lem:laplace}),
\[
c_{L,1}\,s^{-(\theta+2)}
\;\le\;
\int_{[0,\mu_{\mathrm{cut}}]} \mu\,e^{-s\mu}\,dF_b(\mu)
\;\le\;
C_{L,1}\,s^{-(\theta+2)}.
\]
\end{lemma}

\begin{proof}
Split $[0,1/T]\cup(1/T,\mu_{\mathrm{cut}}]$.

\emph{Upper bound.}
On $[0,1/T]$, use $\mu\le 1/T$ and $F_b(1/T)\le C_FT^{-(\theta+1)}$ to get
\[
\int_{[0,1/T]} \mu e^{-s\mu}\,dF_b(\mu) \le \frac{1}{T}F_b(1/T) \lesssim T^{-(\theta+2)}
\asymp s^{-(\theta+2)}.
\]
On $(1/T,\mu_{\mathrm{cut}}]$, use Stieltjes integration by parts with
$f(\mu)=\mu e^{-s\mu}$, so $df(\mu)=(1-s\mu)e^{-s\mu}\,d\mu$:
\[
\int_{(1/T,\mu_{\mathrm{cut}}]} \mu e^{-s\mu}\,dF_b(\mu)
=
\mu_{\mathrm{cut}}e^{-s\mu_{\mathrm{cut}}}F_b(\mu_{\mathrm{cut}})
-\frac{1}{T}e^{-s/T}F_b(1/T)
-\int_{1/T}^{\mu_{\mathrm{cut}}} F_b(\mu)\,(1-s\mu)e^{-s\mu}\,d\mu.
\]
Using $|1-s\mu|\le 1+s\mu$ and $F_b(\mu)\le C_F\mu^{\theta+1}$ yields
\[
\int_{(1/T,\mu_{\mathrm{cut}}]} \mu e^{-s\mu}\,dF_b(\mu)
\le
\mu_{\mathrm{cut}}e^{-s\mu_{\mathrm{cut}}}F_b(\mu_{\mathrm{cut}})
+\frac{1}{T}e^{-s/T}F_b(1/T)
+C_F\int_{1/T}^{\mu_{\mathrm{cut}}} (1+s\mu)e^{-s\mu}\mu^{\theta+1}\,d\mu.
\]
With $t=s\mu$, the last term is
\[
C_F s^{-(\theta+2)}\int_{s/T}^{s\mu_{\mathrm{cut}}} (1+t)e^{-t}t^{\theta+1}\,dt
\le C'_F\,s^{-(\theta+2)}.
\]
For the boundary terms, $ \mu_{\mathrm{cut}}e^{-s\mu_{\mathrm{cut}}}F_b(\mu_{\mathrm{cut}})
\le F_b(\mu_{\mathrm{cut}})\mu_{\mathrm{cut}}e^{-c_1\mu_{\mathrm{cut}}T}=O(e^{-cT})$,
and $(1/T)e^{-s/T}F_b(1/T)=O(T^{-(\theta+2)})\asymp O(s^{-(\theta+2)})$.
This yields the upper bound.

\medskip
\emph{Lower bound.}
Fix any constant $K>1$ such that $c_F K^{\theta+1}-C_F>0$; for instance
\[
K:=2\Big(\frac{C_F}{c_F}\Big)^{\frac{1}{\theta+1}}.
\]
For $T$ large enough we have $K/T\le\mu_{\mathrm{cut}}$, so
$[1/T,K/T]\subset[1/T,\mu_{\mathrm{cut}}]$.
For $\mu\in[1/T,K/T]$ we have $\mu\ge 1/T$ and
$e^{-s\mu}\ge e^{-sK/T}\ge e^{-Kc_2}$ (since $s\le c_2T$). Hence
\[
\int_{[0,\mu_{\mathrm{cut}}]} \mu e^{-s\mu}\,dF_b(\mu)
\ge \int_{(1/T,K/T]} \mu e^{-s\mu}\,dF_b(\mu)
\ge \frac{1}{T}e^{-Kc_2}\big(F_b(K/T)-F_b(1/T)\big).
\]
By Assumption~\ref{ass:wRV},
\[
F_b(K/T)-F_b(1/T)\ge c_F(K/T)^{\theta+1}-C_F(1/T)^{\theta+1}
=(c_FK^{\theta+1}-C_F)\,T^{-(\theta+1)}\gtrsim T^{-(\theta+1)}.
\]
Therefore the last display is $\gtrsim T^{-(\theta+2)}\asymp s^{-(\theta+2)}$ since
$s\asymp T$ on $[c_1T,c_2T]$.
\end{proof}

\begin{lemma}[Tail scaling]\label{lem:ST}
Assume Assumption~\ref{ass:wRV}. Then there exist constants $c_S,C_S>0$ and $T_2\ge 1$
such that for all $T\ge T_2$ and all $\eta\in I$,
\[
c_S(\eta T)^{-(\theta+1)}\le S_T(\eta)\le C_S(\eta T)^{-(\theta+1)}.
\]
In particular, since $\eta\asymp 1$ on $I$, uniformly for $\eta\in I$,
\[
S_T(\eta)\asymp T^{-(\theta+1)}.
\]
\end{lemma}

\begin{proof}
For $\eta\in I$ and $\mu_j\le\mu_{\mathrm{cut}}$, Lemma~\ref{lem:exp-comp} gives
\[
\sum_{\substack{2\le j\le d:\\ \mu_j\le\mu_{\mathrm{cut}}}} b_j^2 e^{-4\eta T\mu_j}
\le S_T(\eta)\le \sum_{\substack{2\le j\le d:\\ \mu_j\le\mu_{\mathrm{cut}}}} b_j^2 e^{-\eta T\mu_j},
\]
so
\[
\mathcal L_{F_b}(4\eta T)\le S_T(\eta)\le \mathcal L_{F_b}(\eta T).
\]
For $\eta\in I$, $s=\eta T$ and $s=4\eta T$ both lie in $[c_1T,c_2T]$ with
$c_1=\eta_0$ and $c_2=8/\mu_1$. Apply Lemma~\ref{lem:laplace} at these two values.
\end{proof}

\begin{lemma}[Tail curvature scaling]\label{lem:tail-curv}
Assume Assumption~\ref{ass:wRV}. Then there exist $c_R,C_R>0$ and $T_R$ such that for all
$T\ge T_R$ and all $\eta\in I$,
\[
c_R(\eta T)^{-(\theta+2)}
\le
\sum_{\substack{2\le j\le d:\\ \mu_j\le\mu_{\mathrm{cut}}}}
b_j^2\,\mu_j\,(1-\eta\mu_j)^{2T-2}
\le
C_R(\eta T)^{-(\theta+2)}.
\]
In particular this sum is $\asymp T^{-(\theta+2)}$ uniformly on $\eta\in I$.
\end{lemma}

\begin{proof}
For $\eta\in I$ and $\mu_j\le\mu_{\mathrm{cut}}$, Lemma~\ref{lem:exp-comp} (with exponent $2T-2$)
gives
\[
\sum b_j^2\mu_j\,e^{-4\eta T\mu_j}
\le
\sum b_j^2\mu_j\,(1-\eta\mu_j)^{2T-2}
\le
\sum b_j^2\mu_j\,e^{-\eta T\mu_j}.
\]
These are exactly $\int \mu e^{-s\mu}\,dF_b(\mu)$ at $s=4\eta T$ and $s=\eta T$.
Apply Lemma~\ref{lem:laplace-mu} to both.
\end{proof}

% ------------------------------------------------------------------
\subsection{Head remainder bounds (uniform on the full interval)}

\begin{lemma}[Uniform head bound from a top-edge gap]\label{lem:head-uniform}
Let $J_{\mathrm{head}}:=\{j\ge 2:\mu_j>\mu_{\mathrm{cut}}\}$. If $J_{\mathrm{head}}=\varnothing$
then $H_T(\eta)\equiv 0$ for all $\eta$. Otherwise, assume Assumption~\ref{ass:gap}.
Then there exists $q\in(0,1)$ depending only on $(\mu_{\mathrm{cut}}/\mu_1,\varepsilon_{\mathrm{gap}})$ such that
for all $\eta\in[1/\mu_1,2/\mu_1]$ and all $\mu\in[\mu_{\mathrm{cut}},\mu_2]$,
\[
|1-\eta\mu|\le q.
\]
Consequently, for all $T\ge 1$ and all $\eta\in[1/\mu_1,2/\mu_1]$,
\[
|H_T(\eta)|\le \|b\|^2 q^{2T-1}.
\]
\end{lemma}

\begin{proof}
If $J_{\mathrm{head}}=\varnothing$, then $H_T(\eta)$ is an empty sum and is identically zero.

Otherwise, fix $\eta\in[1/\mu_1,2/\mu_1]$ and $\mu\in[\mu_{\mathrm{cut}},\mu_2]$. Then
\[
\eta\mu\in\Big[\frac{\mu_{\mathrm{cut}}}{\mu_1},\frac{2\mu_2}{\mu_1}\Big]
\subset \Big(0,2(1-\varepsilon_{\mathrm{gap}})\Big).
\]
For any interval $[a,b]\subset(0,2)$,
$\sup_{x\in[a,b]}|1-x|=\max\{1-a,\ b-1\}$. Hence
\[
|1-\eta\mu|
\le
\max\Big\{1-\frac{\mu_{\mathrm{cut}}}{\mu_1},\ \frac{2\mu_2}{\mu_1}-1\Big\}
\le
\max\Big\{1-\frac{\mu_{\mathrm{cut}}}{\mu_1},\ 1-2\varepsilon_{\mathrm{gap}}\Big\}
=:q<1,
\]
which proves the uniform bound. Summing over $j\in J_{\mathrm{head}}$ gives
\[
|H_T(\eta)|
\le \sum_{j\in J_{\mathrm{head}}} b_j^2 |1-\eta\mu_j|^{2T-1}
\le \|b\|^2 q^{2T-1}.
\]
\end{proof}

% ------------------------------------------------------------------
\subsection{Derivative signs, minimizer location, and balance}

\begin{lemma}[Derivative signs]\label{lem:signs}
Assume Assumptions~\ref{ass:top}, \ref{ass:wRV}, and \ref{ass:gap}. Then there exists $T_3$ such that for all $T\ge T_3$:
\begin{enumerate}[label=(\roman*),leftmargin=*]
\item $L_T'(\eta)<0$ for all $\eta\in(1/\mu_1,\eta_0]$;
\item $L_T'(2/\mu_1)>0$.
\end{enumerate}
\end{lemma}

\begin{proof}
(i) Fix $\eta\in(1/\mu_1,\eta_0]$. Then $1-\eta\mu_1\in[-3/4,0)$, hence
\[
-T b_1^2(1-\eta\mu_1)^{2T-1}=T b_1^2(\eta\mu_1-1)^{2T-1}
\le T b_1^2(3/4)^{2T-1},
\]
which is exponentially small in $T$.

For the tail term, for $\mu_j\le\mu_{\mathrm{cut}}$ we have $1-\eta\mu_j>0$ and each summand
decreases in $\eta$, so $S_T(\eta)\ge S_T(\eta_0)$. By Lemma~\ref{lem:ST}, $S_T(\eta_0)\gtrsim T^{-(\theta+1)}$,
so $T S_T(\eta)\gtrsim T^{-\theta}$.

Finally, by Lemma~\ref{lem:head-uniform}, $T|H_T(\eta)|\le T\|b\|^2 q^{2T-1}$ is exponentially small.
Thus for $T$ large enough, the negative tail term dominates, yielding $L_T'(\eta)<0$.

(ii) At $\eta=2/\mu_1$, the top contribution is
\[
-T b_1^2(1-2)^{2T-1}=T b_1^2.
\]
Lemma~\ref{lem:ST} gives $T S_T(2/\mu_1)=O(T^{-\theta})$, and Lemma~\ref{lem:head-uniform} gives
$T|H_T(2/\mu_1)|\le T\|b\|^2 q^{2T-1}$, both $o(T)$. Hence $L_T'(2/\mu_1)=T b_1^2-o(T)>0$ for $T$ large.
\end{proof}

\begin{lemma}[Existence and localization of a minimizer]\label{lem:minimizer}
For each $T$, $L_T(\eta)$ attains a minimum on $[0,2/\mu_1]$.
Moreover, under Assumptions~\ref{ass:top}, \ref{ass:wRV}, and \ref{ass:gap}, for all $T\ge T_3$,
the (unique) minimizer $\eta_T^\star$ over $[0,2/\mu_1]$ lies in $(\eta_0,2/\mu_1)$
and satisfies $L_T'(\eta_T^\star)=0$.
\end{lemma}

\begin{proof}
Continuity of $L_T$ on the compact interval gives existence.
By Lemma~\ref{lem:signs}(i), $L_T$ is strictly decreasing on $(1/\mu_1,\eta_0]$, so no minimizer lies there.
Also, for $\eta\in[0,1/\mu_1]$ we have $1-\eta\mu_j\ge 0$ for all $j$, so by \eqref{eq:LTprime}
$L_T'(\eta)\le 0$ (and strictly negative unless $b_j=0$ for all $j$), hence no minimizer lies in $[0,1/\mu_1]$.
By Lemma~\ref{lem:signs}(ii), the derivative is positive at $2/\mu_1$, so $2/\mu_1$ is not a minimizer.
Thus $\eta_T^\star\in(\eta_0,2/\mu_1)$ and, being an interior minimizer of a $C^1$ function, satisfies
$L_T'(\eta_T^\star)=0$.
Uniqueness follows from Lemma~\ref{lem:strict-convex}.
\end{proof}

\begin{lemma}[Top--tail balance at the minimizer]\label{lem:balance}
Under Assumptions~\ref{ass:top}, \ref{ass:wRV}, and \ref{ass:gap}, for all $T$ large enough,
\[
b_1^2(\eta_T^\star\mu_1-1)^{2T-1}\asymp S_T(\eta_T^\star),
\]
where the implicit constants depend only on
$\mu_1,\mu_{\mathrm{cut}},\varepsilon_{\mathrm{gap}},\theta,c_F,C_F,|b_1|,\|b\|$.
\end{lemma}

\begin{proof}
At the interior minimizer, $L_T'(\eta_T^\star)=0$, i.e.
\[
b_1^2(\eta_T^\star\mu_1-1)^{2T-1} = S_T(\eta_T^\star)+H_T(\eta_T^\star).
\]
On $\eta_T^\star\in(\eta_0,2/\mu_1)$, Lemma~\ref{lem:ST} yields $S_T(\eta_T^\star)\gtrsim T^{-(\theta+1)}$,
while Lemma~\ref{lem:head-uniform} implies $|H_T(\eta_T^\star)|\lesssim q^{2T}$ for a fixed $q<1$.
Thus $H_T(\eta_T^\star)$ is exponentially smaller than $S_T(\eta_T^\star)$ and can be absorbed into constants.
\end{proof}

% ------------------------------------------------------------------
\subsection{Main deterministic expansion theorem}

\begin{theorem}[Deterministic optimal step size via weighted spectral regularity]\label{thm:det}
Assume Assumptions~\ref{ass:top}, \ref{ass:wRV}, and \ref{ass:gap}. Let $\eta_T^\star$ be the minimizer
of $L_T(\eta;H,b)$ over $\eta\in(0,2/\mu_1)$. Then as $T\to\infty$ (along any regime where
Assumption~\ref{ass:wRV} holds),
\[
\eta_T^\star
=
\frac{2}{\mu_1}
-\frac{(\theta+1)\log T}{2\mu_1 T}
+O\!\left(\frac{1}{T}\right),
\]
where the $O(1/T)$ constant depends only on
\[
\mu_1,\ \mu_{\mathrm{cut}},\ \varepsilon_{\mathrm{gap}},\ \theta,\ c_F,\ C_F,\ |b_1|,\ \|b\|.
\]
\end{theorem}

\begin{proof}
By Lemma~\ref{lem:minimizer}, $\eta_T^\star\in(\eta_0,2/\mu_1)$ and $L_T'(\eta_T^\star)=0$.
Write $\eta_T^\star=\frac{2}{\mu_1}-\delta_T$ with $\delta_T\in(0,1/(4\mu_1))$,
so $\eta_T^\star\mu_1-1=1-\mu_1\delta_T$ with $\mu_1\delta_T\in(0,1/4)$.

By Lemma~\ref{lem:balance} and Lemma~\ref{lem:ST},
\[
b_1^2(1-\mu_1\delta_T)^{2T-1}\asymp S_T(\eta_T^\star)\asymp T^{-(\theta+1)}.
\]
Taking logs gives
\begin{equation}\label{eq:logstep}
(2T-1)\log(1-\mu_1\delta_T)=-(\theta+1)\log T+O(1).
\end{equation}
For $x\in(0,1/4)$ one has $-x-x^2\le \log(1-x)\le -x$. Applying this with $x=\mu_1\delta_T$
in \eqref{eq:logstep} yields the two-sided bound
\[
\frac{(\theta+1)\log T-C_1}{2\mu_1 T}\le \delta_T\le \frac{(\theta+1)\log T+C_2}{2\mu_1 T}
\]
for constants $C_1,C_2$ depending only on the fixed parameters listed in the theorem.
Therefore
\[
\delta_T=\frac{(\theta+1)\log T}{2\mu_1 T}+O\!\left(\frac{1}{T}\right),
\]
which implies the claimed expansion for $\eta_T^\star$.
\end{proof}

\begin{theorem}[Curvature at the optimal step size]\label{thm:det-curv}
Assume Assumptions~\ref{ass:top}, \ref{ass:wRV}, and \ref{ass:gap}. Let $\eta_T^\star$ be the minimizer in
Theorem~\ref{thm:det}. Then there exist constants $0<c_{\mathrm{curv}}\le C_{\mathrm{curv}}<\infty$
(depending only on $\mu_1,\mu_{\mathrm{cut}},\varepsilon_{\mathrm{gap}},\theta,c_F,C_F,|b_1|,\|b\|$) such that for all
$T$ large enough,
\[
c_{\mathrm{curv}}\,\mu_1\,T^{1-\theta}
\;\le\;
L_T''(\eta_T^\star)
\;\le\;
C_{\mathrm{curv}}\,\mu_1\,T^{1-\theta}.
\]
In particular $L_T''(\eta_T^\star)>0$ for all sufficiently large $T$.
\end{theorem}

\begin{proof}
Decompose \eqref{eq:LTsecond} into top, tail, and head:
\[
L_T''(\eta)=L''_{\mathrm{top}}(\eta)+L''_{\mathrm{tail}}(\eta)+L''_{\mathrm{head}}(\eta),
\]
where $j=1$ is top, $\{j\ge2:\mu_j\le \mu_{\mathrm{cut}}\}$ is tail, and $\{j\ge2:\mu_j>\mu_{\mathrm{cut}}\}$ is head.

\emph{Lower bound.}
Since $L_T''\ge L''_{\mathrm{top}}$ and $\eta_T^\star\in(\eta_0,2/\mu_1)$,
\[
L_T''(\eta_T^\star)
\ge
T(2T-1)\,b_1^2\mu_1\,(\eta_T^\star\mu_1-1)^{2T-2}.
\]
By Lemma~\ref{lem:balance},
$b_1^2(\eta_T^\star\mu_1-1)^{2T-1}\asymp S_T(\eta_T^\star)$, hence
$(\eta_T^\star\mu_1-1)^{2T-2}\asymp S_T(\eta_T^\star)/b_1^2$ up to a bounded factor.
Therefore
\[
L_T''(\eta_T^\star)\gtrsim \mu_1\,T^2\,S_T(\eta_T^\star).
\]
Finally Lemma~\ref{lem:ST} gives $S_T(\eta_T^\star)\gtrsim T^{-(\theta+1)}$, so
$L_T''(\eta_T^\star)\gtrsim \mu_1 T^{1-\theta}$.

\emph{Upper bound.}
For the top term, Lemma~\ref{lem:balance} and Lemma~\ref{lem:ST} give
$(\eta_T^\star\mu_1-1)^{2T-2}\lesssim T^{-(\theta+1)}$, hence
\[
L''_{\mathrm{top}}(\eta_T^\star)\lesssim \mu_1 T^2\cdot T^{-(\theta+1)}=\mu_1 T^{1-\theta}.
\]
For the tail term, Lemma~\ref{lem:tail-curv} yields
\[
L''_{\mathrm{tail}}(\eta_T^\star)
\lesssim
T^2\cdot T^{-(\theta+2)}=T^{-\theta}\le \mu_1 T^{1-\theta}
\]
for all large $T$. For the head term, Lemma~\ref{lem:head-uniform} gives $|1-\eta\mu|\le q<1$ uniformly, hence
\[
0\le L''_{\mathrm{head}}(\eta_T^\star)
\le
T(2T-1)\sum_{\mu_j>\mu_{\mathrm{cut}}} b_j^2\mu_j\,q^{2T-2}
\le
T(2T-1)\mu_2\|b\|^2 q^{2T-2},
\]
which is exponentially small in $T$. Summing the three bounds gives the upper bound.
\end{proof}

\begin{remark}[What the curvature scaling says]
Theorem~\ref{thm:det-curv} gives $L_T''(\eta_T^\star)\asymp \mu_1 T^{1-\theta}$.
Thus curvature stays bounded away from $0$ (indeed, diverges) if $\theta<1$, is order $1$ if
$\theta=1$, and decays polynomially if $\theta>1$.
\end{remark}

% ------------------------------------------------------------------
\subsection{Uniform curvature on a shrinking window around $\eta_T^\star$}

\begin{lemma}[Uniform curvature near $\eta_T^\star$ on a $1/T$ window]
\label{lem:curv-window}
Assume Assumptions~\ref{ass:top}, \ref{ass:wRV}, and \ref{ass:gap}. Fix any $M>0$ and
define the shrinking window
\[
W_{T,M}
:=
\Big[\eta_T^\star-\frac{M}{T},\ \eta_T^\star+\frac{M}{T}\Big]\cap I,
\qquad
I=\Big[\eta_0,\frac{2}{\mu_1}\Big].
\]
Then there exist constants $0<c_M\le C_M<\infty$ and $T_M\ge 1$
(depending only on $\mu_1,\mu_{\mathrm{cut}},\varepsilon_{\mathrm{gap}},\theta,c_F,C_F,|b_1|,\|b\|$ and $M$)
such that for all $T\ge T_M$,
\[
c_M\,\mu_1\,T^{1-\theta}
\ \le\
L_T''(\eta;H,b)
\ \le\
C_M\,\mu_1\,T^{1-\theta},
\qquad \forall \eta\in W_{T,M}.
\]
In particular, if $\theta\ge 1$ then $\sup_{\eta\in W_{T,M}} L_T''(\eta)=O(1)$ as $T\to\infty$.
\end{lemma}

\begin{proof}
Decompose \eqref{eq:LTsecond} into top, tail, and head contributions:
\[
L_T''(\eta)=L_{\mathrm{top}}''(\eta)+L_{\mathrm{tail}}''(\eta)+L_{\mathrm{head}}''(\eta),
\]
where
\[
L_{\mathrm{top}}''(\eta)=T(2T-1)b_1^2\mu_1(1-\eta\mu_1)^{2T-2},
\]
and $L_{\mathrm{tail}}''$ (resp.\ $L_{\mathrm{head}}''$) is the sum over
$\{j\ge 2:\mu_j\le \mu_{\mathrm{cut}}\}$ (resp.\ $\{j\ge 2:\mu_j>\mu_{\mathrm{cut}}\}$).

\medskip\noindent
\emph{Step 1: the window stays away from the boundary.}
Write $\eta_T^\star=\frac{2}{\mu_1}-\delta_T$ with $\delta_T>0$.
By Theorem~\ref{thm:det}, $\delta_T=\frac{(\theta+1)\log T}{2\mu_1 T}+O(1/T)$, hence
$\delta_T T\to\infty$. Therefore, for any fixed $M>0$ and all $T$ sufficiently large,
$\eta_T^\star+\frac{M}{T}\le \frac{2}{\mu_1}$ and $W_{T,M}\subset I$ is nonempty.

Define
\[
a(\eta):=\eta\mu_1-1\in(0,1] \qquad \text{for }\eta\in\Big(\frac{1}{\mu_1},\frac{2}{\mu_1}\Big].
\]
For $\eta\in W_{T,M}$ we have $a(\eta)=1-\mu_1\big(\frac{2}{\mu_1}-\eta\big)
=1-\mu_1\big(\delta_T+O(1/T)\big)\to 1$.
Hence, for all $T$ large enough (depending on $M$), we may assume
\begin{equation}\label{eq:a-lower}
\inf_{\eta\in W_{T,M}} a(\eta)\ge \frac12.
\end{equation}

\medskip\noindent
\emph{Step 2: the top curvature is stable across $W_{T,M}$.}
For $\eta\in W_{T,M}$, use the mean value theorem for $\log a(\eta)$:
\[
\big|\log a(\eta)-\log a(\eta_T^\star)\big|
\le \sup_{\xi\in W_{T,M}}\frac{1}{a(\xi)}\,|a(\eta)-a(\eta_T^\star)|.
\]
By \eqref{eq:a-lower} and $a(\eta)=\mu_1\eta-1$,
\[
\big|\log a(\eta)-\log a(\eta_T^\star)\big|
\le 2\,\mu_1\,|\eta-\eta_T^\star|
\le \frac{2\mu_1 M}{T}.
\]
Multiplying by $2T-2\le 2T$ and exponentiating yields
\[
e^{-4\mu_1 M}
\le
\frac{a(\eta)^{2T-2}}{a(\eta_T^\star)^{2T-2}}
\le
e^{4\mu_1 M},
\qquad \forall \eta\in W_{T,M}.
\]
Consequently,
\begin{equation}\label{eq:top-stable}
e^{-4\mu_1 M}\,L_{\mathrm{top}}''(\eta_T^\star)
\ \le\
L_{\mathrm{top}}''(\eta)
\ \le\
e^{4\mu_1 M}\,L_{\mathrm{top}}''(\eta_T^\star),
\qquad \forall \eta\in W_{T,M}.
\end{equation}

\medskip\noindent
\emph{Step 3: tail and head curvature are uniformly negligible on $I$.}
By Lemma~\ref{lem:tail-curv}, uniformly for $\eta\in I$,
\[
0\le L_{\mathrm{tail}}''(\eta)
=
T(2T-1)\sum_{\substack{2\le j\le d:\\ \mu_j\le\mu_{\mathrm{cut}}}}
b_j^2\mu_j(1-\eta\mu_j)^{2T-2}
\ \lesssim\
T^2\cdot T^{-(\theta+2)}
=
T^{-\theta}.
\]
By Lemma~\ref{lem:head-uniform}, there exists $q\in(0,1)$ such that uniformly for $\eta\in I$,
\[
0\le L_{\mathrm{head}}''(\eta)
\le
T(2T-1)\sum_{\substack{2\le j\le d:\\ \mu_j>\mu_{\mathrm{cut}}}}
b_j^2\mu_j\,q^{2T-2}
\ \lesssim\ T^2 q^{2T-2}.
\]

\medskip\noindent
\emph{Step 4: comparison to the curvature scale at $\eta_T^\star$.}
By Theorem~\ref{thm:det-curv}, there exist constants $0<c_{\mathrm{curv}}\le C_{\mathrm{curv}}<\infty$
such that for all large $T$,
\begin{equation}\label{eq:curv-star-scale}
c_{\mathrm{curv}}\mu_1 T^{1-\theta}
\le
L_T''(\eta_T^\star)
\le
C_{\mathrm{curv}}\mu_1 T^{1-\theta}.
\end{equation}
Moreover, since $L_T''(\eta_T^\star)=L_{\mathrm{top}}''(\eta_T^\star)+L_{\mathrm{tail}}''(\eta_T^\star)+L_{\mathrm{head}}''(\eta_T^\star)$
and the last two terms are $O(T^{-\theta})+O(T^2q^{2T})$, we also have
$L_{\mathrm{top}}''(\eta_T^\star)\asymp \mu_1 T^{1-\theta}$ with constants depending only on the same fixed parameters.

Combine \eqref{eq:top-stable} with the bounds on $L_{\mathrm{tail}}''$ and $L_{\mathrm{head}}''$ to obtain, uniformly for $\eta\in W_{T,M}$,
\[
L_T''(\eta)
=
L_{\mathrm{top}}''(\eta)+O(T^{-\theta})+O(T^2q^{2T})
\asymp
\mu_1 T^{1-\theta}.
\]
This yields the desired two-sided bounds with constants $c_M,C_M$.
\end{proof}

\begin{remark}[Quadratic control of $L_T$ near $\eta_T^\star$]
\label{rem:quad-window}
Under the assumptions of Lemma~\ref{lem:curv-window}, for $T$ large enough and all $\eta\in W_{T,M}$,
Taylor's theorem with integral remainder and the fact that $L_T'(\eta_T^\star)=0$ imply
\[
\frac{c_M}{2}\,\mu_1\,T^{1-\theta}\,(\eta-\eta_T^\star)^2
\ \le\
L_T(\eta)-L_T(\eta_T^\star)
\ \le\
\frac{C_M}{2}\,\mu_1\,T^{1-\theta}\,(\eta-\eta_T^\star)^2.
\]
\end{remark}


% ------------------------------------------------------------------
\section{Gaussian-sketched random features: model and concentration inputs}
\label{sec:rf}

\subsection{Random features and the induced quadratic problem}

Let $\mathcal H$ be a real separable Hilbert space. Let $\psi\in\mathcal H$ be centered with
trace-class covariance operator
\[
\Sigma:=\mathbb E[\psi\otimes\psi],
\]
assumed self-adjoint, positive, and trace-class, with eigenpairs $(\lambda_k,u_k)$
and a spectral gap $\gamma:=\lambda_1-\lambda_2>0$.

\paragraph{Gaussian sketch (basis-free).}
Let $g_1,\dots,g_N$ be i.i.d.\ standard Gaussian elements on $\mathcal H$ (isonormal process),
so that for any $f\in\mathcal H$, $\langle g_i,f\rangle\sim\mathcal N(0,\|f\|^2)$ and the
$g_i$ are independent. Define the sketched features $\phi\in\mathbb R^N$ by
\[
\phi_i := \frac{1}{\sqrt N}\,\langle g_i,\psi\rangle,\qquad i=1,\dots,N.
\]
Define also the Hilbert-valued Gaussian vectors
\[
X_i := \Sigma^{1/2}g_i\in\mathcal H,\qquad i=1,\dots,N,
\]
and the linear map $Z:\mathbb R^N\to\mathcal H$ by $Z e_i = X_i$ (equivalently $Z=[X_1,\dots,X_N]$).
Then the RF covariance matrix and RF linear term are
\[
H_N:=\mathbb E[\phi\phi^\top]=\frac1N Z^\top Z,
\qquad
b_N:=\mathbb E[y\phi]=\frac{1}{\sqrt N}\,Z^\top \Sigma^{1/2}\bar w,
\]
for noiseless targets $y=\langle \bar w,\psi\rangle$.
Thus the RF step-size problem is exactly the deterministic problem with $(H,b)=(H_N,b_N)$
in dimension $d=N$.

\subsection{Almost sure SPD of $H_N$}

\begin{lemma}[Almost sure SPD of $H_N$]\label{lem:SPD}
Assume $\Sigma$ has at least $N$ strictly positive eigenvalues (e.g.\ $\lambda_k>0$ for all $k$).
Then $H_N$ is SPD almost surely, hence $v_N^\star=H_N^{-1}b_N$ is well-defined almost surely.
\end{lemma}

\begin{proof}
We have $H_N=\frac1N Z^\top Z$ with $Z=[X_1,\dots,X_N]$ and $X_i=\Sigma^{1/2}g_i$ i.i.d.\ Gaussian in $\mathcal H$
with covariance $\Sigma$. If $\mathrm{rank}(\Sigma)\ge N$, then the law of $(X_1,\dots,X_N)$ is nondegenerate on an
$N$-dimensional injective subspace, which implies the columns of $Z$ are linearly independent almost surely.
Thus $Z^\top Z$ is SPD almost surely.
\end{proof}

\subsection{Covariance-operator concentration (Koltchinskii--Lounici)}

Define the sample covariance operator
\[
\widehat\Sigma_N:=\frac1N\sum_{i=1}^N X_i\otimes X_i = \frac1N Z Z^\top.
\]
The nonzero eigenvalues of $\widehat\Sigma_N$ coincide with those of $H_N=\frac1N Z^\top Z$
(with the same multiplicities), and in particular the top eigenvalues match.

\begin{lemma}[Koltchinskii--Lounici concentration; informal form]\label{lem:KL}
Let $X$ be centered Gaussian in a separable Hilbert space with covariance $\Sigma$ and
$\widehat\Sigma_N$ its sample covariance from $N$ i.i.d.\ samples. Then
\[
\mathbb E\|\widehat\Sigma_N-\Sigma\|_{\mathrm{op}}
\asymp
\|\Sigma\|_{\mathrm{op}}\Big(\sqrt{\frac{r(\Sigma)}{N}}\ \vee\ \frac{r(\Sigma)}{N}\Big),
\]
and for all $t\ge 1$, with probability at least $1-e^{-t}$,
\[
\Big|\|\widehat\Sigma_N-\Sigma\|_{\mathrm{op}}-\mathbb E\|\widehat\Sigma_N-\Sigma\|_{\mathrm{op}}\Big|
\lesssim
\|\Sigma\|_{\mathrm{op}}\Big(\sqrt{\frac{t}{N}}\ \vee\ \frac{t}{N}\Big),
\]
where $r(\Sigma)$ is an intrinsic dimension parameter (effective rank in the Hilbert case).
See Koltchinskii--Lounici~\cite{KL2017} for precise definitions and constants.
\end{lemma}

\subsection{Top eigenvalue perturbation and a uniform relative gap}

Let $\widehat\lambda_1^{(N)}\ge \widehat\lambda_2^{(N)}\ge\cdots$ denote the eigenvalues of $H_N$
(equivalently of $\widehat\Sigma_N$).

\begin{lemma}[High-probability top-eigenvalue control]\label{lem:top-eig}
Let $\gamma=\lambda_1-\lambda_2>0$. There exist $c>0$ and $N_0$ such that for all $N\ge N_0$,
\[
\mathbb P\big(\|\widehat\Sigma_N-\Sigma\|_{\mathrm{op}}>\gamma/2\big)\le e^{-cN}.
\]
On the event $\{\|\widehat\Sigma_N-\Sigma\|_{\mathrm{op}}\le\gamma/2\}$ one has
$|\widehat\lambda_1^{(N)}-\lambda_1|\le \gamma/2$ and $\widehat\lambda_1^{(N)}\in[\lambda_1/2,3\lambda_1/2]$.
\end{lemma}

\begin{proof}
By Lemma~\ref{lem:KL}, $\mathbb E\|\widehat\Sigma_N-\Sigma\|_{\mathrm{op}}=O(N^{-1/2})$.
Choose $N_0$ so that $\mathbb E\|\widehat\Sigma_N-\Sigma\|_{\mathrm{op}}\le\gamma/4$ for $N\ge N_0$.
Using the high-probability part of Lemma~\ref{lem:KL} with $t=cN$ and choosing $c>0$ sufficiently small
(depending on $\gamma/\|\Sigma\|_{\mathrm{op}}$) makes the fluctuation term at most $\gamma/4$, yielding the claim.
Weyl's inequality gives $|\widehat\lambda_1^{(N)}-\lambda_1|\le\|\widehat\Sigma_N-\Sigma\|_{\mathrm{op}}$.
\end{proof}

\begin{lemma}[Uniform relative top gap on the same good event]\label{lem:rf-gap}
On the event $\{\|\widehat\Sigma_N-\Sigma\|_{\mathrm{op}}\le \gamma/4\}$,
\[
\frac{\widehat\lambda_2^{(N)}}{\widehat\lambda_1^{(N)}}
\le
\frac{\lambda_1-3\gamma/4}{\lambda_1-\gamma/4}
=1-\varepsilon_{\mathrm{RF}},
\qquad
\varepsilon_{\mathrm{RF}}:=\frac{\gamma/2}{\lambda_1-\gamma/4}>0.
\]
\end{lemma}

\begin{proof}
Weyl gives $\widehat\lambda_1^{(N)}\ge \lambda_1-\gamma/4$ and
$\widehat\lambda_2^{(N)}\le \lambda_2+\gamma/4=\lambda_1-3\gamma/4$. Divide the two bounds.
\end{proof}

% ------------------------------------------------------------------
\subsection{Population capacity and source assumptions}\label{subsec:cap-source}

Let $(\lambda_k,u_k)_{k\ge 1}$ be the eigenpairs of $\Sigma$, ordered so that
$\lambda_1>\lambda_2\ge\lambda_3\ge\cdots>0$.

\begin{assumption}[Capacity]\label{ass:cap}
There exist $\alpha>1$ and constants $c_\lambda,C_\lambda>0$ such that
\[
c_\lambda k^{-\alpha} \le \lambda_k \le C_\lambda k^{-\alpha}
\qquad \forall k\ge 1.
\]
\end{assumption}

\begin{assumption}[Source (including the top mode)]\label{ass:source}
There exist $\beta>0$ and $g\in\mathcal H$ such that $\bar w = \Sigma^\beta g$.
In the eigenbasis $(u_k)$, write $g_k:=\langle g,u_k\rangle$ and assume there are constants
$0<c_g\le C_g<\infty$ such that
\[
c_g \le |g_k|^2 \le C_g
\qquad \forall k\ge 1.
\]
\end{assumption}

Let $b_\infty:=\Sigma\bar w$ and $b_{\infty,k}:=\langle b_\infty,u_k\rangle$.
Fix $\mu_{\mathrm{cut}}\le \lambda_1/8$ and define the population weighted CDF (away from the top mode)
\[
F_\infty(x)
:=\sum_{\substack{k\ge 2\\ \lambda_k\le \min\{x,\mu_{\mathrm{cut}}\}}}
b_{\infty,k}^2.
\]

\begin{proposition}[Population weighted regular variation from capacity and source]\label{prop:pop-wRV}
Under Assumptions~\ref{ass:cap} and \ref{ass:source}, there exist constants
$c_\infty,C_\infty>0$ and $x_0>0$ such that for all $x\in(0,x_0]$,
\[
c_\infty x^{\theta+1}\le F_\infty(x)\le C_\infty x^{\theta+1},
\qquad
\theta:=2\beta+1-\frac{1}{\alpha}>0.
\]
\end{proposition}

\begin{proof}
In the eigenbasis of $\Sigma$, $\bar w_k=\lambda_k^\beta g_k$, hence
$b_{\infty,k}=\lambda_k\bar w_k=\lambda_k^{\beta+1}g_k$ and $b_{\infty,k}^2\asymp \lambda_k^{2(\beta+1)}$.
With $\lambda_k\asymp k^{-\alpha}$, the condition $\lambda_k\le x$ is equivalent to $k\gtrsim x^{-1/\alpha}$, giving
\[
F_\infty(x)\asymp \sum_{k\gtrsim x^{-1/\alpha}} k^{-2\alpha(\beta+1)}
\asymp x^{2(\beta+1)-1/\alpha}=x^{\theta+1}.
\]
\end{proof}

\subsection{Top eigenvalue perturbation and inverse moments}

Let $\widehat\lambda_1^{(N)}$ denote the top eigenvalue of $H_N$ (equivalently of $\widehat\Sigma_N$).

\begin{lemma}[High-probability top-eigenvalue control]\label{lem:top-eig}
Let $\gamma=\lambda_1-\lambda_2>0$. There exists $c>0$ and $N_0$ such that for all $N\ge N_0$,
\[
\mathbb P\big(\|\widehat\Sigma_N-\Sigma\|_{\mathrm{op}}>\gamma/2\big)\le e^{-cN}.
\]
On the event $\{\|\widehat\Sigma_N-\Sigma\|_{\mathrm{op}}\le\gamma/2\}$ one has
$|\widehat\lambda_1^{(N)}-\lambda_1|\le \gamma/2$ and in particular
$\widehat\lambda_1^{(N)}\in[\lambda_1/2,3\lambda_1/2]$.
\end{lemma}

\begin{proof}
By Lemma~\ref{lem:KL}, $\mathbb E\|\widehat\Sigma_N-\Sigma\|_{\mathrm{op}}=O(N^{-1/2})$.
Choose $N_0$ so that $\mathbb E\|\widehat\Sigma_N-\Sigma\|_{\mathrm{op}}\le\gamma/4$ for $N\ge N_0$.
Then take $t=cN$ and use the high-probability part of Lemma~\ref{lem:KL} to ensure the fluctuation term is $\le\gamma/4$,
yielding $\mathbb P(\|\widehat\Sigma_N-\Sigma\|_{\mathrm{op}}>\gamma/2)\le e^{-cN}$.
Weyl's inequality gives $|\widehat\lambda_1^{(N)}-\lambda_1|\le\|\widehat\Sigma_N-\Sigma\|_{\mathrm{op}}$.
\end{proof}

\begin{lemma}[Uniform inverse-moment bounds]\label{lem:inv-mom-bdd}
For all $N>4$,
\[
\mathbb E\!\left[\frac{1}{\widehat\lambda_1^{(N)}}\right]\le \frac{1}{\lambda_1}\frac{N}{N-2},
\qquad
\mathbb E\!\left[\frac{1}{(\widehat\lambda_1^{(N)})^2}\right]\le \frac{1}{\lambda_1^2}\frac{N^2}{(N-2)(N-4)}.
\]
In particular, $\mathbb E[1/(\widehat\lambda_1^{(N)})^2]=O(1)$ uniformly in $N$.
\end{lemma}

\begin{proof}
Decompose $\Sigma=\lambda_1 u_1u_1^\top+\Sigma_\perp$ with $\Sigma_\perp\succeq 0$. Then
\[
H_N=A^\top\Sigma A\succeq \lambda_1 A^\top u_1u_1^\top A = \lambda_1 z z^\top,\qquad z:=A^\top u_1.
\]
Hence $\widehat\lambda_1^{(N)}\ge \lambda_1\|z\|_2^2$. Since $A_{ki}\sim\mathcal N(0,1/N)$ and $\|u_1\|=1$,
$z\sim \mathcal N(0,(1/N)I_N)$ and $\|z\|_2^2\stackrel{d}{=}\chi^2_N/N$. Therefore
\[
\frac{1}{\widehat\lambda_1^{(N)}}\le \frac{1}{\lambda_1}\frac{N}{\chi^2_N},\qquad
\frac{1}{(\widehat\lambda_1^{(N)})^2}\le \frac{1}{\lambda_1^2}\frac{N^2}{(\chi^2_N)^2}.
\]
Using $\mathbb E[1/\chi^2_N]=1/(N-2)$ and $\mathbb E[1/(\chi^2_N)^2]=1/((N-2)(N-4))$ for $N>4$ gives the result.
\end{proof}

\begin{lemma}[Inverse-moment expansions to order $1/N$]\label{lem:inv-exp}
Assume $\gamma:=\lambda_1-\lambda_2>0$. Then
\[
\mathbb E\!\left[\frac{1}{\widehat\lambda_1^{(N)}}\right]=\frac{1}{\lambda_1}+O\!\left(\frac{1}{N}\right),
\qquad
\mathbb E\!\left[\frac{2}{\widehat\lambda_1^{(N)}}\right]=\frac{2}{\lambda_1}+O\!\left(\frac{1}{N}\right),
\]
where the $O(1/N)$ constants depend only on $\lambda_1,\gamma,\|\Sigma\|_{\mathrm{op}}$ and the
intrinsic-dimension parameter $r(\Sigma)$ appearing in Lemma~\ref{lem:KL}.
\end{lemma}

\begin{proof}
This is the same perturbation-plus-Taylor-expansion argument as in the draft: write $E=\widehat\Sigma_N-\Sigma$,
work on the good event $\|E\|_{\mathrm{op}}\le \gamma/2$, use eigenvalue perturbation to get
$\widehat\lambda_1^{(N)}-\lambda_1=\langle u_1,Eu_1\rangle+O(\|E\|_{\mathrm{op}}^2/\gamma)$,
then Taylor expand $x\mapsto 1/x$ around $\lambda_1$ and use $\mathbb E\|E\|_{\mathrm{op}}^2=O(1/N)$ from KL.
The contribution of the complement event is exponentially small and controlled via Lemma~\ref{lem:inv-mom-bdd}.
\end{proof}

% ------------------------------------------------------------------
\subsection{RF weighted regularity on scales $\mu\sim 1/T$ (imported local-law input)}

Fix $\mu_{\mathrm{cut}}\le \lambda_1/8$. For $H_N$, define the RF weighted CDF away from the top mode:
\[
F^{(N)}_{b_N}(x):=\sum_{\substack{2\le j\le N:\\ \widehat\lambda_j^{(N)}\le \min\{x,\mu_{\mathrm{cut}}\}}}
\langle b_N,\widehat v_j^{(N)}\rangle^2,
\]
where $(\widehat\lambda_j^{(N)},\widehat v_j^{(N)})$ are the eigenpairs of $H_N$.

\begin{assumption}[RF weighted spectral regularity on relevant scales]\label{ass:RF-wRV}
Fix $\mu_{\mathrm{cut}}\le \lambda_1/8$ and $\theta>0$. There exist constants
$c_F,C_F,c_0,C_0>0$, $T_0\ge 1$, and $c_T>0$ such that for all $N$ large enough and all integers
$T$ in the regime
\[
T_0\le T\le c_T N^\alpha,
\]
there is an event $\mathcal E_{N,T}$ with \emph{overwhelming probability}
(i.e.\ for every $D>0$, $\mathbb P(\mathcal E_{N,T}^c)\le N^{-D}$ for $N$ large enough) on which:
\begin{enumerate}[label=(\alph*),leftmargin=*]
\item (\emph{Top eigenvalue and gap control})
$\widehat\lambda_1^{(N)}\in[\lambda_1/2,2\lambda_1]$ and
$\widehat\lambda_2^{(N)}/\widehat\lambda_1^{(N)}\le 1-\varepsilon_{\mathrm{RF}}$
for a fixed $\varepsilon_{\mathrm{RF}}>0$ depending only on $(\lambda_1,\gamma)$.
\item (\emph{Top-mode overlap and norm control})
$|\langle b_N,\widehat v_1^{(N)}\rangle|\ge c_0$ and $\|b_N\|_2\le C_0$.
\item (\emph{Weighted regular variation on $[1/T,\mu_{\mathrm{cut}}]$})
For all $x\in[1/T,\mu_{\mathrm{cut}}]$,
\[
c_F x^{\theta+1}\le F^{(N)}_{b_N}(x)\le C_F x^{\theta+1}.
\]
\end{enumerate}
\end{assumption}

\begin{remark}[Why the regime $T\le c_TN^\alpha$ appears]\label{rem:T-regime}
Under capacity $\lambda_k\asymp k^{-\alpha}$ one expects $\mu_{\min}(H_N)\asymp N^{-\alpha}$.
Then $T\le c_TN^\alpha$ ensures $1/T\gtrsim \mu_{\min}(H_N)$ so the window $[1/T,\mu_{\mathrm{cut}}]$ lies above the
spectral floor (see Remark~\ref{rem:finite-d}).
\end{remark}

% ------------------------------------------------------------------
\section{RF optimal step size (and curvature): final theorems}

Let $\eta^\star(T,N)$ denote a minimizer of the RF population excess loss after $T$ GD
steps with step size $\eta$ on the RF quadratic, i.e.\ a minimizer of $L_T(\eta;H_N,b_N)$
over stable $\eta\in(0,2/\widehat\lambda_1^{(N)})$.

\begin{theorem}[Gaussian RF: optimal constant step size (high probability and in expectation)]
\label{thm:RF-clean}
Let $\Sigma$ be a trace-class, self-adjoint, positive operator on a separable Hilbert space
with eigenvalues $\lambda_1>\lambda_2\ge\cdots>0$ and gap $\gamma:=\lambda_1-\lambda_2>0$.
Let $\psi$ be centered with covariance $\Sigma$, and let $g_1,\dots,g_N$ be i.i.d.\ standard
Gaussians on the Hilbert space (isonormal process). Define
\[
X_i:=\Sigma^{1/2}g_i,\qquad Z:=[X_1,\dots,X_N],\qquad
H_N:=\frac1N Z^\top Z,\qquad b_N:=\frac{1}{\sqrt N}Z^\top\Sigma^{1/2}\bar w.
\]
Consider population GD on the quadratic defined by $(H_N,b_N)$ for $T$ steps, and let
$\eta^\star(T,N)$ be a minimizer of $L_T(\eta;H_N,b_N)$ over stable $\eta\in(0,2/\widehat\lambda_1^{(N)})$,
where $\widehat\lambda_1^{(N)}$ is the top eigenvalue of $H_N$.

Assume capacity/source:
\begin{enumerate}[label=(\roman*),leftmargin=*]
\item (Capacity) $\lambda_k\asymp k^{-\alpha}$ for some $\alpha>1$.
\item (Source) $\bar w=\Sigma^\beta g$ with some $\beta>0$, and in the eigenbasis $(u_k)$,
      $|g_k|^2\asymp 1$ for all $k\ge 1$ (in particular $g_1\neq 0$).
\end{enumerate}
Let $\theta:=2\beta+1-\frac{1}{\alpha}>0$ and fix $\mu_{\mathrm{cut}}\le \lambda_1/8$.
Fix a horizon regime $T_0\le T\le c_T N^\alpha$.

Assume furthermore that there exists an event $\mathcal E_{N,T}$ on which the following hold:
\begin{enumerate}[label=(\alph*),leftmargin=*]
\item $\widehat\lambda_1^{(N)}\in[\lambda_1/2,2\lambda_1]$ and
      $\widehat\lambda_2^{(N)}/\widehat\lambda_1^{(N)}\le 1-\varepsilon_{\mathrm{RF}}$
      for some fixed $\varepsilon_{\mathrm{RF}}>0$ depending only on $(\lambda_1,\gamma)$;
\item $|\langle b_N,\widehat v_1^{(N)}\rangle|\ge c_0$ and $\|b_N\|_2\le C_0$;
\item (RF weighted regularity on $[1/T,\mu_{\mathrm{cut}}]$) for all $x\in[1/T,\mu_{\mathrm{cut}}]$,
\[
c_F x^{\theta+1}\le
F^{(N)}_{b_N}(x)
:=
\sum_{\substack{2\le j\le N:\\ \widehat\lambda_j^{(N)}\le \min\{x,\mu_{\mathrm{cut}}\}}}
\langle b_N,\widehat v_j^{(N)}\rangle^2
\le C_F x^{\theta+1}.
\]
\end{enumerate}
and that $\mathbb P(\mathcal E_{N,T}^c)\le \delta_{N,T}$.

Then there exist deterministic constants $C_T<\infty$ and $N_0$ (depending only on
$\lambda_1,\gamma,\mu_{\mathrm{cut}},\alpha,\beta$ and the constants in (a)--(c)) such that
for all $N\ge N_0$ and all $T_0\le T\le c_TN^\alpha$:

\medskip
\noindent\textbf{(1) High-probability expansion.}
On $\mathcal E_{N,T}$,
\[
\eta^\star(T,N)
=
\frac{2}{\widehat\lambda_1^{(N)}}
-\frac{(\theta+1)\log T}{2\widehat\lambda_1^{(N)}\,T}
+O\!\left(\frac{1}{T}\right),
\]
where the $O(1/T)$ constant is bounded by $C_T$ uniformly over $\mathcal E_{N,T}$.

\medskip
\noindent\textbf{(2) Expectation expansion with explicit failure-probability remainder.}
Moreover,
\[
\mathbb E[\eta^\star(T,N)]
=
\frac{2}{\lambda_1}
-\frac{(\theta+1)\log T}{2\lambda_1 T}
+O\!\left(\frac{1}{T}\right)
+O\!\left(\frac{1}{N}\right)
+O\!\big(\delta_{N,T}^{1/2}\big).
\]
In particular, if $\delta_{N,T}\le N^{-D}$ for some $D>2$ uniformly over
$T\le c_TN^\alpha$, then the last term is $o(1/N)$ and can be absorbed into $O(1/N)$.
\end{theorem}


\begin{thebibliography}{99}

\bibitem{KL2017}
V.~Koltchinskii and K.~Lounici.
\newblock Concentration inequalities and moment bounds for sample covariance operators.
\newblock \emph{Bernoulli} 23(1):110--133, 2017. Also available as arXiv:1405.2468.

\bibitem{Chouard2022}
C.~Chouard.
\newblock Quantitative deterministic equivalent of sample covariance matrices with a general dependence structure.
\newblock arXiv:2211.13044, 2022.

\bibitem{DLM2024}
L.~Defilippis, B.~Loureiro, and T.~Misiakiewicz.
\newblock Dimension-free deterministic equivalents and scaling laws for random feature regression.
\newblock arXiv:2405.15699, 2024. Also in NeurIPS 2024.

\bibitem{Bordelon2023}
B.~Bordelon, A.~Atanasov, and C.~Pehlevan.
\newblock A dynamical model of neural scaling laws.
\newblock \emph{Advances in Neural Information Processing Systems}, 2023.

\bibitem{Karkada2024}
D.~Karkada, J.~Turnbull, Y.~Liu, and J.~B.~Simon.
\newblock Predicting kernel regression learning curves from only raw data statistics.
\newblock arXiv:2402.09371, 2024.

\bibitem{Simon2023}
J.~B. Simon et al.
\newblock More is Better in Modern Machine Learning.
\newblock arXiv:2311.14646, 2023.

\end{thebibliography}

\end{document}
